\chapter*{Resumen}
\addcontentsline{toc}{chapter}{Resumen}

Existe una literatura amplia que predice hospitalización y mortalidad por
COVID-19 con datos mexicanos, y su desempeño puede considerarse un problema
resuelto. Esta tesis desplaza el objeto de estudio: no pregunta si un modelo
acierta, sino \emph{si acierta por igual}, con énfasis en el sesgo geográfico
entre entidades federativas, una dimensión asociada a la desigualdad estructural
del sistema de salud en México y prácticamente inexplorada.

Sobre la base abierta de COVID-19 de la Dirección General de Epidemiología
(DGE), restringida a casos confirmados del año 2021
($n=\cohorteN$; mortalidad observada \cohorteTasaDef\,\%), se entrena un modelo
de referencia de mortalidad al primer contacto, con control estricto de fuga de
información, y se calibra mediante regresión isotónica sobre un bloque de
calibración independiente. El modelo alcanza un AUROC de \aurocGlobal{} y un
puntaje de Brier de \brierCalibrado{} tras calibrar (frente a \brierCrudo{} sin
calibrar). Ese modelo se somete después a una auditoría de equidad por entidad,
sexo y grupo etario, y a las tres familias de mitigación de sesgo
--pre, in y post-procesamiento-- sobre un punto de operación clínico común de
\sensObjetivo\,\% de sensibilidad.

Los resultados confirman inequidad geográfica sustantiva: la sensibilidad varía
\gapTPRentidad{} entre entidades y el AUROC \gapAUROCentidad{}, con el peor
desempeño en \entPeor{} (AUROC \entPeorAUROC) y el mejor en \entMejor{}
(\entMejorAUROC). La brecha persiste al restringir el análisis a adultos de 60
años o más, lo que descarta que se explique solo por composición etaria. En
cambio, \emph{no} se explica por la marginación de la entidad: el índice de
CONAPO 2020 no muestra asociación con el desempeño ($r=\corrMargAurocr$,
$p=\corrMargAurocp$) ni aparece gradiente al agregar por grado de marginación.
Entre
las intervenciones evaluadas, únicamente el post-procesamiento por umbrales
específicos de entidad reduce la disparidad de forma apreciable --de \gapBase{}
a \gapPost{}, una reducción de \reduccionPost\,\%-- y lo hace sin costo de
desempeño global. La reponderación previa al entrenamiento y la restricción de
\emph{equalized odds} durante el entrenamiento resultan insuficientes o
inestables con 32 grupos de tamaño muy desigual. Documentar esa asimetría entre
familias de mitigación es, en sí, un resultado.

Dos análisis de robustez acotan la lectura. Al excluir por completo una entidad
del entrenamiento y evaluarla, el desempeño apenas cambia (pérdida mediana de
AUROC \loeoMediana) y las brechas no se acentúan: la inequidad no proviene de la
cobertura de datos. Y al restringir la cohorte a los pacientes hospitalizados
--la población en la que el desenlace es posible, pues fallecer implica estar
registrado como hospitalizado-- el AUROC cae a \hospAUROC{} y la disparidad se
desplaza del umbral hacia el ordenamiento, lo que sugiere que la mitigación por
umbrales sería insuficiente en la población clínicamente crítica.

La contribución es la primera auditoría de equidad geográfica de un modelo
clínico sobre datos abiertos mexicanos, acompañada de una frontera de Pareto
equidad--desempeño utilizable como instrumento de decisión antes del despliegue.

\vspace{1em}
\noindent\textbf{Palabras clave:} equidad algorítmica, sesgo geográfico,
estratificación de riesgo, COVID-19, datos abiertos de salud, calibración,
\emph{equalized odds}, frontera de Pareto.

\cleardoublepage
